\documentclass[a4paper,12pt,twoside]{memoir}

% Castellano
\usepackage[spanish,es-tabla]{babel}
\selectlanguage{spanish}
\usepackage[utf8]{inputenc}
\usepackage[T1]{fontenc}
\usepackage{lmodern} % scalable font
\usepackage{microtype}
\usepackage{placeins}

\RequirePackage{booktabs}
\RequirePackage[table]{xcolor}
\RequirePackage{xtab}
\RequirePackage{multirow}

% Links
\usepackage[colorlinks]{hyperref}
\hypersetup{
	allcolors = {red}
}

% Ecuaciones
\usepackage{amsmath}

% Rutas de fichero / paquete
\newcommand{\ruta}[1]{{\sffamily #1}}

% Párrafos
\nonzeroparskip


% Imagenes
\usepackage{graphicx}
\newcommand{\imagen}[2]{
	\begin{figure}[!h]
		\centering
		\includegraphics[width=0.9\textwidth]{#1}
		\caption{#2}\label{fig:#1}
	\end{figure}
	\FloatBarrier
}

\newcommand{\imagenflotante}[2]{
	\begin{figure}%[!h]
		\centering
		\includegraphics[width=0.9\textwidth]{#1}
		\caption{#2}\label{fig:#1}
	\end{figure}
}



% El comando \figura nos permite insertar figuras comodamente, y utilizando
% siempre el mismo formato. Los parametros son:
% 1 -> Porcentaje del ancho de página que ocupará la figura (de 0 a 1)
% 2 --> Fichero de la imagen
% 3 --> Texto a pie de imagen
% 4 --> Etiqueta (label) para referencias
% 5 --> Opciones que queramos pasarle al \includegraphics
% 6 --> Opciones de posicionamiento a pasarle a \begin{figure}
\newcommand{\figuraConPosicion}[6]{%
  \setlength{\anchoFloat}{#1\textwidth}%
  \addtolength{\anchoFloat}{-4\fboxsep}%
  \setlength{\anchoFigura}{\anchoFloat}%
  \begin{figure}[#6]
    \begin{center}%
      \Ovalbox{%
        \begin{minipage}{\anchoFloat}%
          \begin{center}%
            \includegraphics[width=\anchoFigura,#5]{#2}%
            \caption{#3}%
            \label{#4}%
          \end{center}%
        \end{minipage}
      }%
    \end{center}%
  \end{figure}%
}

%
% Comando para incluir imágenes en formato apaisado (sin marco).
\newcommand{\figuraApaisadaSinMarco}[5]{%
  \begin{figure}%
    \begin{center}%
    \includegraphics[angle=90,height=#1\textheight,#5]{#2}%
    \caption{#3}%
    \label{#4}%
    \end{center}%
  \end{figure}%
}
% Para las tablas
\newcommand{\otoprule}{\midrule [\heavyrulewidth]}
%
% Nuevo comando para tablas pequeñas (menos de una página).
\newcommand{\tablaSmall}[5]{%
 \begin{table}
  \begin{center}
   \rowcolors {2}{gray!35}{}
   \begin{tabular}{#2}
    \toprule
    #4
    \otoprule
    #5
    \bottomrule
   \end{tabular}
   \caption{#1}
   \label{tabla:#3}
  \end{center}
 \end{table}
}

%
%Para el float H de tablaSmallSinColores
\usepackage{float}

%
% Nuevo comando para tablas pequeñas (menos de una página).
\newcommand{\tablaSmallSinColores}[5]{%
 \begin{table}[H]
  \begin{center}
   \begin{tabular}{#2}
    \toprule
    #4
    \otoprule
    #5
    \bottomrule
   \end{tabular}
   \caption{#1}
   \label{tabla:#3}
  \end{center}
 \end{table}
}

\newcommand{\tablaApaisadaSmall}[5]{%
\begin{landscape}
  \begin{table}
   \begin{center}
    \rowcolors {2}{gray!35}{}
    \begin{tabular}{#2}
     \toprule
     #4
     \otoprule
     #5
     \bottomrule
    \end{tabular}
    \caption{#1}
    \label{tabla:#3}
   \end{center}
  \end{table}
\end{landscape}
}

%
% Nuevo comando para tablas grandes con cabecera y filas alternas coloreadas en gris.
\newcommand{\tabla}[6]{%
  \begin{center}
    \tablefirsthead{
      \toprule
      #5
      \otoprule
    }
    \tablehead{
      \multicolumn{#3}{l}{\small\sl continúa desde la página anterior}\\
      \toprule
      #5
      \otoprule
    }
    \tabletail{
      \hline
      \multicolumn{#3}{r}{\small\sl continúa en la página siguiente}\\
    }
    \tablelasttail{
      \hline
    }
    \bottomcaption{#1}
    \rowcolors {2}{gray!35}{}
    \begin{xtabular}{#2}
      #6
      \bottomrule
    \end{xtabular}
    \label{tabla:#4}
  \end{center}
}

%
% Nuevo comando para tablas grandes con cabecera.
\newcommand{\tablaSinColores}[6]{%
  \begin{center}
    \tablefirsthead{
      \toprule
      #5
      \otoprule
    }
    \tablehead{
      \multicolumn{#3}{l}{\small\sl continúa desde la página anterior}\\
      \toprule
      #5
      \otoprule
    }
    \tabletail{
      \hline
      \multicolumn{#3}{r}{\small\sl continúa en la página siguiente}\\
    }
    \tablelasttail{
      \hline
    }
    \bottomcaption{#1}
    \begin{xtabular}{#2}
      #6
      \bottomrule
    \end{xtabular}
    \label{tabla:#4}
  \end{center}
}

%
% Nuevo comando para tablas grandes sin cabecera.
\newcommand{\tablaSinCabecera}[5]{%
  \begin{center}
    \tablefirsthead{
      \toprule
    }
    \tablehead{
      \multicolumn{#3}{l}{\small\sl continúa desde la página anterior}\\
      \hline
    }
    \tabletail{
      \hline
      \multicolumn{#3}{r}{\small\sl continúa en la página siguiente}\\
    }
    \tablelasttail{
      \hline
    }
    \bottomcaption{#1}
  \begin{xtabular}{#2}
    #5
   \bottomrule
  \end{xtabular}
  \label{tabla:#4}
  \end{center}
}



\definecolor{cgoLight}{HTML}{EEEEEE}
\definecolor{cgoExtralight}{HTML}{FFFFFF}

%
% Nuevo comando para tablas grandes sin cabecera.
\newcommand{\tablaSinCabeceraConBandas}[5]{%
  \begin{center}
    \tablefirsthead{
      \toprule
    }
    \tablehead{
      \multicolumn{#3}{l}{\small\sl continúa desde la página anterior}\\
      \hline
    }
    \tabletail{
      \hline
      \multicolumn{#3}{r}{\small\sl continúa en la página siguiente}\\
    }
    \tablelasttail{
      \hline
    }
    \bottomcaption{#1}
    \rowcolors[]{1}{cgoExtralight}{cgoLight}

  \begin{xtabular}{#2}
    #5
   \bottomrule
  \end{xtabular}
  \label{tabla:#4}
  \end{center}
}




\graphicspath{ {./img/} }

% Capítulos
\chapterstyle{bianchi}
\newcommand{\capitulo}[2]{
	\setcounter{chapter}{#1}
	\setcounter{section}{0}
	\chapter*{#2}
	\addcontentsline{toc}{chapter}{#2}
	\markboth{#2}{#2}
}

% Apéndices
\renewcommand{\appendixname}{Apéndice}
\renewcommand*\cftappendixname{\appendixname}

\newcommand{\apendice}[1]{
	%\renewcommand{\thechapter}{A}
	\chapter{#1}
}

\renewcommand*\cftappendixname{\appendixname\ }

% Formato de portada
\makeatletter
\usepackage{xcolor}
\newcommand{\tutor}[1]{\def\@tutor{#1}}
\newcommand{\course}[1]{\def\@course{#1}}
\definecolor{cpardoBox}{HTML}{E6E6FF}
\def\maketitle{
  \null
  \thispagestyle{empty}
  % Cabecera ----------------
\noindent\includegraphics[width=\textwidth]{cabecera}\vspace{1cm}%
  \vfill
  % Título proyecto y escudo informática ----------------
  \colorbox{cpardoBox}{%
    \begin{minipage}{.8\textwidth}
      \vspace{.5cm}\Large
      \begin{center}
      \textbf{TFG del Grado en Ingeniería Informática}\vspace{.6cm}\\
      \textbf{\LARGE\@title{}}
      \end{center}
      \vspace{.2cm}
    \end{minipage}

  }%
  \hfill\begin{minipage}{.20\textwidth}
    \includegraphics[width=\textwidth]{escudoInfor}
  \end{minipage}
  \vfill
  % Datos de alumno, curso y tutores ------------------
  \begin{center}%
  {%
    \noindent\LARGE
    Presentado por \@author{}\\ 
    en Universidad de Burgos --- \@date{}\\
    Tutor: \@tutor{}\\
  }%
  \end{center}%
  \null
  \cleardoublepage
  }
\makeatother


% Datos de portada
\title{IG-Study \\Documentación Técnica}
\author{Miguel Pérez López de Echazarreta}
\tutor{Dr. Jose Manuel Galán Ordax  \\ y Dr. José Ignacio Santos Martín}
\date{\today}

\begin{document}

\maketitle



\cleardoublepage



%%%%%%%%%%%%%%%%%%%%%%%%%%%%%%%%%%%%%%%%%%%%%%%%%%%%%%%%%%%%%%%%%%%%%%%%%%%%%%%%%%%%%%%%



\frontmatter


\clearpage

% Indices
\tableofcontents

\clearpage

\listoffigures

\clearpage

\listoftables

\clearpage

\mainmatter

\appendix

\apendice{Plan de Proyecto Software}



\section{Introducción}


En esta sección vamos a tratar la planificación seguida durante el proyecto. La planificación es un una parte fundamental a la hora de embarcarse en un proyecto, y es ahí cuando se toman se fijan temas como:
\begin{itemize}
    \item Los objetivos
    \item El tiempo
    \item La financiación
\end{itemize}
Por este motivo la planificación se va a dividir en:
\begin{itemize}
    \item Planificación temporal
    \item Estudio de viabilidad
\end{itemize}




\section{Planificación temporal}

Para la planificación temporal, haremos uso de una metodología ágil: SCRUM.
Se sigue una estrategia con el desarrollo incremental. La duración de los Sprints era de una o dos semanas, dependiendo de las tutorías futuras posibles y factores externos que pudieran entorpecer el ritmo del proyecto.
Se describirán ahora los Sprints realizados:

\subsection{Sprint 0 (6/3/19 – 12/3/19):}

En la primera reunión de planificación del proyecto se hablo de las ideas que iban a llevarse a cabo. Se mandaron las primeras tareas a realizar.
Las tareas que realizar fueron:
\begin{itemize}
    \item Descarga de ciertos Softwares a usar durante 1el proyecto, tales como Balsamic o ZenHub.
    \item Estudio de la políticas y licencias de Instagram.
    \item Investigación del funcionamiento de la API.
\end{itemize}

\subsection{Sprint 1 (13/3/19 – 20/3/19):}

Las tareas a realizar en este Sprint fueron:
\begin{itemize}
    \item Crear una cuenta de Instagram
    \item Conectarme a la API
\end{itemize}


\subsection{Sprint 2 (21/3/19 – 27/3/19):}

Las tareas a realizar en este Sprint fueron:
\begin{itemize}
    \item Investigación sobre como hacer el reconocimiento de la gente de Burgos
    \item Hacer una búsqueda de palabras que se usen en las descripciones de los perfiles.
\end{itemize}


\subsection{Sprint 3 (28/3/19 – 10/4/19):}

Las tareas a realizar en este Sprint fueron:
\begin{itemize}
    \item Espera a los permisos de la conexión de la base de datos
    \item Búsqueda de Bases de Datos
    \item Conexión VPN
\end{itemize}

\subsection{Sprint 5 (11/4/19 – 17/4/19):}

Las tareas a realizar en este Sprint fueron:
\begin{itemize}
    \item Preparar servidor en ordenador personal
    \item Descarga de Screen para Terminal
    \item Instalación de bibliotecas
\end{itemize}

\subsection{Sprint 6 (18/4/19 – 1/5/19):}

Las tareas a realizar en este Sprint fueron:
\begin{itemize}
    \item Instalación del conector de la base de datos
    \item Recopilación de datos
    \item Extracción y maquetación de los datos
\end{itemize}

\subsection{Sprint 7 (2/5/19 – 15/5/19):}

Las tareas a realizar en este Sprint fueron:
\begin{itemize}
    \item Instalación de la biblioteca de Louvain
    \item División por comunidades de los datos
    \item Investigar si hay alguna comunidad clave
\end{itemize}

\subsection{Sprint 8 (9/5/19 – 22/5/19):}

Las tareas a realizar en este Sprint fueron:
\begin{itemize}
    \item Estudio de personas relevantes global
    \item Estudio por comunidades de los usuarios más relevantes
\end{itemize}

\subsection{Sprint 9 (23/5/19 – 5/6/19):}

Las tareas a realizar en este Sprint fueron:
\begin{itemize}
    \item Estudio de Hashtags más usados en toda la red
    \item Estudio de Hashtags más relevantes por comunidades
\end{itemize}

\subsection{Sprint 10 (6/6/19 – 19/6/19):}

Las tareas a realizar en este Sprint fueron:
\begin{itemize}
    \item Aprendizaje de Flask
    \item Realización de la página web donde mostrar los resultados
\end{itemize}

\subsection{Sprint 11 (20/6/19 – 3/7/19):}

Las tareas a realizar en este Sprint fueron:
\begin{itemize}
    \item Pulir detalles del proyecto
    \item Terminar documentación 
    \item Entregar proyecto
\end{itemize}

\subsection{Sprint 12 (4/7/19 - 10/8/19):}

Las tareas a realizar en este Sprint fueron:
\begin{itemize}
    \item Presentar el proyecto
\end{itemize}



\section{Estudio de viabilidad}

Este apartado muestra el estudio realizado sobre la viabilidad económica y legal del producto a entregar.

\subsection{Viabilidad económica}

En la siguiente sección se realiza una aproximación de los costes que tendría que asumir una empresa para llevar adelante este proyecto.

\subsubsection{Coste del personal}
El transcurso del proyecto ha sido de 3 meses desde que el desarrollador junior se puso a trabajar en él. Suponiendo que el desarrollador tiene un salario neto de 1100 euros:

\begin{table}[!h]
	\centering
	\begin{tabular}{@{}l|l@{}}
		\toprule
		Concepto & Coste  \\
		\midrule
		Salario mensual neto & 1100 euros \\
		Seguridad Social (29,9 \%) & 596,91 euros \\
		Salario mensual bruto & 1.696,91 euros \\
		\midrule
		Total 3 meses & 5.090,73 euros \\
		\bottomrule
	\end{tabular}
	\caption{Coste del personal}
	\label{tab:personal}
\end{table}

La retribución a la Seguridad Social es de un 29,9\% siguiendo los siguientes puntos a valorar:


\begin{itemize}
	\item 23,6\% de contingencias comunes.
	\item 5,5\% de desempleo de tipo general.
	\item 0,6\% de formación profesional.
	\item 0,2\% del fondo de garantía salarial.
\end{itemize}

\subsubsection{Coste del material}
En el coste material habría que mirar tanto el coste del Hardware como del Software. En mi caso a nivel de software no ha habido ningún coste puesto que he ido buscando siempre usar software libre. 
Por este motivo solo hay que valorar el coste material del Hardware.
Para la realización del proyecto se ha utilizado el ordenador personal del alumno que tiene un valor de 500 euros al cual se le estima un tiempo para amortizarlo de 6 años. Dado que el tiempo que ha sido de 3 meses:


\begin{table}[!h]
	\centering
	\begin{tabular}{l|l|l}
		\toprule
		Concepto & Coste & Coste amortizado  \\
		\midrule
		Ordenador & 500 euros & 20,83 euros \\
		\bottomrule
	\end{tabular}
	\caption{Coste del Hardware}
	\label{tab:material}
\end{table}

\subsubsection{Coste Total}
La suma de todos los costes es:

\begin{table}[!h]
	\centering
	\begin{tabular}{@{}l|l@{}}
		\toprule
		Concepto & Coste  \\
		\midrule
		Hardware & 20,83 \\
		Personal & 5.090,73 euros \\
		\midrule
		Total & 5.111,56 euros \\
		\bottomrule
	\end{tabular}
	\caption{Coste Total del Proyecto}
	\label{tab:costeTotal}
\end{table}

\subsubsection{Beneficios}

La forma de obtener beneficios es vender la aplicación para realizar estudios de mercado. Basándome en un precio teórico de 250 euros por cada estudio de mercado que se requiera, se tardaría en amortizar el producto 25 estudios de mercado. A partir de ese momento habría beneficios.


\subsection{Viabilidad legal}

En este siguiente apartado se va a estudiar la viabilidad legal del proyecto. Hay que analizar todas las licencias de las bibliotecas que se han usado para ver que licencia se podría aplicar a nuestro producto software.

\begin{table}[!h]
	\centering
	\begin{tabular}{@{}l|l@{}}
		\toprule
		Librería & Licencia  \\
		\midrule
		Python and modules integrated & PSFL[Zotero] \\
		Instagram-API-python & MIT License[Zotero] \\
		Networkx & 3-clause BSD license [Zotero] \\
		\bottomrule
	\end{tabular}
	\caption{Tabla librerías}
	\label{tab:librerias}
\end{table}

Como podemos observar en la tablas, las licencias de las bibliotecas utilizadas son bastante permisivas y tienden a favorecer el open source.
Este proyecto lo publico con la licencia “Apache 2.0”, que dictamina lo siguiente:


\begin{table}[!h]
	\centering
	\begin{tabular}{l|l|l}
		\toprule
		Permissions & Limitations & Conditions  \\
		\midrule
		Commercial use & Trademark & License and copyright notice \\
		Modification & Liability & State changes \\
		Distribution & Warranty \\
		Patent use \\
		Private use \\
		\bottomrule
	\end{tabular}
	\caption{Tabla Licencias}
	\label{tab:licencia}
\end{table}
\apendice{Especificación de Requisitos}

\section{Introducción}
Este anexo va a describir y formalizar los requisitos que debe de cumplir nuestra herramienta software en cuanto a funcionalidad se respecta.
\section{Objetivos generales}
Los objetivos generales del producto son:
\begin{itemize}
    \item Obtención de la red
    \item Estudio de los usuarios
    \item Implementación de una web para mostrar resultados
\end{itemize}



\section{Catalogo de requisitos}
Los requisitos se dividen en dos tipos:
\subsection{Requisitos Funcionales:}
\begin{itemize}
    \item RF1. Crear perfil en Instagram
    \item RF2. Usar el perfil para loguearse
    \item RF3. Crear red de usuarios a estudiar
    \item RF4. Guardar datos en la Base de Datos
    \item RF5. Realizar división por comunidades
    \item RF6. Búsqueda de Usuarios relevantes
\end{itemize}

\subsection{Requisitos no Funcionales:}
\begin{itemize}
    \item RNF1. Modularidad: El software puede ser continuado por otro alumno, por lo cual debe de ser modular.
    \item RNF2. Visibilidad: La web que enseña los datos debe de mostrar bien los datos, de manera intuitiva y explicativa
\end{itemize}


\subsection{Catálogo de requisitos:}

\subsubsection{Caso de uso 1:}
\begin{itemize}
\item Crear perfil en Instagram
\item Descripción: El usuario debe de crear un perfil en Instagram
\item Requisitos: Ninguno
\item Precondiciones: Tener un correo electrónico y conexión a internet
\item Secuencia normal 	Paso	 Acción
\begin{itemize}
    \item 1	Entrar en la página de Instagram 
    \item 2	Darse de alta con un correo electrónico
\end{itemize}
\item Postcondiciones: Se obtiene una manera de continuar en Instagram haciendo uso de su API
\item Excepciones: Tener un usuario previamente
\item Importancia: Alta
\item Urgencia: Alta
\item Comentarios:Ninguno
\end{itemize}

\subsubsection{Caso de uso 2:}
\begin{itemize}
\item Usar perfil de IG para loguearse
\item Descripción: El usuario debe insertar el usuario y contraseña en los scripts para poder hacer uso del proyecto
\item Requisitos: RF1
\item Precondiciones: Ninguna
\item Secuencia normal 	Paso	 Acción
\begin{itemize}
    \item 1	Entrar en los scripts 
    \item 2	Escribir el usuario y contraseña
\end{itemize}
\item Postcondiciones: Se puede hacer uso de los distintos procesos implementados
\item Excepciones: Ninguna
\item Importancia: Alta
\item Urgencia: Alta
\item Comentarios:Ninguno
\end{itemize}

\subsubsection{Caso de uso 3:}
\begin{itemize}
\item Creación de la red de usuarios
\item Descripción: El usuario debe insertar los valores que van a buscarse en las descripciones para buscar personas que lo cumplan
\item Requisitos: RF2
\item Precondiciones: Tener agregados previamente a 5 personas que vayan a estar en la red para tener un punto para comenzar la búsqueda.
\item Secuencia normal 	Paso	 Acción
\begin{itemize}
    \item 1	Entrar en los scripts 
    \item 2	Escribir la lista de términos que van a buscarse
    \item 3	Correr el Script durante un tiempo largo
\end{itemize}
\item Postcondiciones: Se seguirá a las personas que cumplen los requisitos con nuestra cuenta de Instagram
\item Excepciones: Ninguna
\item Importancia: Alta
\item Urgencia: Alta
\item Comentarios: Puede entrarse en la cuenta de Instagram para ver cuántas personas se han agregado para poder parar el script cuando sean suficientes. Puede durar días este paso.
\end{itemize}


\subsubsection{Caso de uso 4:}
\begin{itemize}
\item Guardar datos en la base de datos
\item Descripción: El usuario debe de correr la aplicación encargada de ver las publicaciones de la gente de la red obteniendo la información.
\item Requisitos: RF2, RF3
\item Precondiciones: Tener instalado Mysql
\item Secuencia normal 	Paso	 Acción
\begin{itemize}
    \item 1	Ejecutar el script
\end{itemize}
\item Postcondiciones: Ninguna
\item Excepciones: Tener un usuario previamente
\item Importancia: Alta
\item Urgencia: Alta
\item Comentarios: El tiempo requerido va a ser proporcional al número de publicaciones de los usuarios y número de usuarios en sí.
\end{itemize}


\subsubsection{Caso de uso 5:}
\begin{itemize}
\item Dividir por comunidades
\item Descripción: El usuario debe de correr la aplicación encargada de marcar las comunidades encontradas en la red.
\item Requisitos: RF4
\item Precondiciones: Tener los datos en CSV
\item Secuencia normal 	Paso	 Acción
\begin{itemize}
    \item 1	Ejecutar el script teniendo los datos en la misma ubicación del programa
\end{itemize}
\item Postcondiciones: Se obtiene una manera de continuar en Instagram haciendo uso de su API
\item Excepciones: Tener un usuario previamente
\item Importancia: Alta
\item Urgencia: Alta
\item Comentarios:Ninguno
\end{itemize}


\subsubsection{Caso de uso 6:}
\begin{itemize}
\item Búsqueda de Usuarios relevantes 
\item Descripción: Se obtendrá una lista con los usuarios más relevantes de la red y comunidades
\item Requisitos: RF2, RF3
\item Precondiciones: Usar Google Colab
\item Secuencia normal 	Paso	 Acción
\begin{itemize}
    \item 1	Subir el notebook a google colab
    \item 2	Ejecutar el Notebook y seguir los pasos que nos va indicando
\end{itemize}
\item Postcondiciones: Ninguna
\item Excepciones: Ninguna
\item Importancia: Alta
\item Urgencia: Alta
\item Comentarios:Ninguno
\end{itemize}



\apendice{Especificación de diseño}

\section{Introducción}
Este anexo va a describir la razón de ser del procedimiento que he seguido a la hora de organizar el proyecto. También se describirán las razones por las cuales se han tomado así las decisiones.
\section{Diseño de datos}
Ahora se explicarán como están implementados los datos en la base de datos en MySQL.
\subsection{Tablas de la base de datos}
Para guardar la información requerida de la API de Instagram usamos la base de datos con las relaciones de la siguiente manera:
\imagen{diagramaDB.png}{Estructura de la Base de Datos}



\begin{itemize}
    \item Persona: Almacena las personas que van a formar la red.
    \item Persona\_sigue: Almacena la relación de seguimiento entre usuarios.
    \item Publicación: Almacena la publicación con la persona que la ha realizado. No se almacena la foto, solo la fecha y el texto a investigar.
    \item Publicación\_Hashtag: Almacena la publicación con los diferentes hashtags que contenga.
    \item Publicación\_like: Almacena la publicación con la gente de la red que haya interaccionando con ella dando un “me gusta”.
\end{itemize}


\section{Diseño procedimental}
\imagen{diagrama-secuencia.png}{Diseño procedimental del diagrama de Secuencia}
Como se puede observar, es un modelo secuencial que sigue siempre la misma dirección. Se ha intentado mantener una sencillez que ayude tanto al usuario como al programador que tenga que hacer uso de ella.
A continuación, comienzo la explicación del proceso que se observa del diagrama:


\begin{itemize}
    \item 1 Se introducen los datos, tanto del usuario, como para la búsqueda y creación de la red.
    \item 2 Se procede a la creación de la red.
    \item 3 Se extraen los datos necesarios de la API y se almacenan en la base de datos de MySQL.
    \item 4 Se analizan los datos globales de la base de datos.
    \item 5 Se calculan las comunidades y se guardan los datos como csv.
    \item 6 Se analizan los datos para obtener los resultados de las distintas comunidades.
    \item 7 Se muestran los resultados en la web.
\end{itemize}


\section{Diseño arquitectónico}
El código del programa se ha llevado a cabo dividiendo el programa en 3 partes:

\begin{itemize}
    \item Scripts principales: En los scripts principales están las principales funciones de creación de red, extracción de datos y análisis de los datos.
\end{itemize}
\imagen{scriptPython.PNG}{Estructura de los Scripts de Python}


\begin{itemize}
    \item Static: en este apartado están los estáticos que se van a usar en la web. Imágenes, datos en csv, fuentes, etc.
\end{itemize}

\imagen{static.png}{Estructura del directorio donde almacenar imagenes, csv, fuentes}

\begin{itemize}
    \item Templates: en este apartado se encuentra los archivos en HTML para poder mostrar la web.
\end{itemize}

\imagen{templates.PNG}{Estructura de los templates de la web}
\apendice{Documentación técnica de programación}

\section{Introducción}
Este anexo va a describir la documentación referente a la programación. Se incluirá también la instalación necesaria de todos los componentes para poder seguir desarrollando el proyecto.
\section{Estructura de directorios}
A continuación, se van a mostrar los directorios de este proyecto que están en el directorio raíz.

\imagen{directorios.png}{Estructura de los directorios del proyecto}


\section{Manual del programador}
En este apartado se facilitarán los programas y librerías que son necesarios para seguir con el desarrollo del proyecto.
Para trabajar con este proyecto son necesarios los siguientes programas y librerias:
\begin{itemize}
    \item Pycharm
    \item Python
    \item Pip
    \item Instagram-API-python
    \item Networkx
    \item Python-Louvain
    \item Flask
\end{itemize}


\subsection{Instalación de PyCharm}
Pycharm ha sido el IDE que se eligió para programar los scripts de nuestro proyecto. Para la instalación hay que descargarse el software desde su página web \cite{pycharm} eligiendo la versión que queramos instalar.
Gracias a la cuenta de la Universidad de Burgos, con la cuenta universitaria podemos descargarnos la versión profesional de manera gratuita, aunque la versión gratuita también puede usarse.
\imagen{pycharm.png}{Modelo que vamos a elegir de PyCharm}

Cuando esté descargado, solo es seguir el instalador que no requiere de ningún conocimiento técnico, solo seguir los pasos.

\subsection{Python}
Python es el lenguaje principal sobre el que se programa todo el proyecto. Esto requiere que sea imprescindible que se tenga instalado para poder seguir con el proyecto. Se ha utilizado la versión 2.7.15. En el momento de la entrega del proyecto hay versiones más avanzadas, sin embargo, esta cumplía con todos los requisitos necesarios para el uso de las librerías.
Para instalarlo hay que ir a la pagina oficial de Python\cite{python} elegir el sistema operativo con el que se va a trabajar y la versión que queramos utilizar. Una vez descarguemos el asistente de instalación, solo habría que ejecutarlo y seguir los pasos que se indiquen.

\subsection{Pip}
Pip es una herramienta que proporciona Python que ayuda sobremanera en la instalación, actualización y desinstalación de bibliotecas desde línea de comandos. Las librerías que se van a usar van a ser descargadas usando pip, por lo tanto, en caso de que no se haya instalado de manera automática con Python, es muy recomendable descargarlo.

\subsection{Instagram-API-Python}\cite{git-API-python}
Instagram-API-Python es la biblioteca que nos acerca la api de Instagram a Python y va a ser usado para la creación de la red y extracción de la información. Es requerida para esos 2 primeros pasos y se instala escribiendo en la línea de comandos:

\begin{itemize}
    \item pip install InstagramApi
\end{itemize}

para su uso en los script de Python hay que importar escribiendo:

\begin{itemize}
    \item from InstagramAPI import InstagramAPI
\end{itemize}

\subsection{Networkx}\cite{networkx}
Instagram-API-Python es la biblioteca que nos acerca la api de Instagram a Python y va a ser usado para la creación de la red y extracción de la información. Es requerida para esos 2 primeros pasos y se instala escribiendo en la línea de comandos:

\begin{itemize}
    \item pip install networkx
\end{itemize}

\subsection{Python-Louvain}\cite{python-louvain}
Es una biblioteca que funciona como adjunta de Networkx y que nos permite hacer la división clustering. Podría hacerse división usando otro método, pero como este es el que hemos elegido, se requiere de esta biblioteca paralela.
Para su instalación basta con escribir en la línea de comandos:

\begin{itemize}
    \item pip install python-louvain
\end{itemize}

\subsection{Flask}
Flask es el framework que se va a utilizar para la web que se encarga de mostrar los resultados finales. En caso de que se quiera hacer algún cambio o correrlo hay que tenerlo instalado. Su instalación es sencilla. Basta con escribir en la línea de comandos:

\begin{itemize}
    \item pip install Flask
\end{itemize}


\section{Compilación, instalación y ejecución del proyecto}
Para la ejecución del proyecto es necesario descargarse el proyecto del repositorio de Github, donde está alojado.Para acceder al repositorio solo hay que acceder al siguiente enlace (ZOTERO). A continuación, pulsaremos el botón verde de la parte de la derecha que dice “Clone or download”
Salvo que tengamos la opción de clonar el git, particularmente me parece más simple descargar el zip.
\subsection{Instalación}
Este proyecto no requiere instalación más allá de tener todas las librerías necesarias instaladas previamente para su uso.

\subsection{Ejecución}
Para la ejecución he dividido las distintas fases en distintos directorios.
Hay que ejecutar cada archivo de Python por orden, pero previamente hay que introducir el usuario de Instagram que se vaya a usar, el termino de búsqueda de las descripciones y la base de datos.
Para la extracción de los datos del mysql a CSV hay que escribir como administrador el siguiente comando en mysql:
\begin{itemize}
    \item select IDseguidor, IDseguido from persona\_sigue into outfile '/var/lib/mysql-files/prueba.csv' fields enclosed by '/' terminated by ';' escaped by '"' LINES terminated by '/r/n';
\end{itemize}

Hay que tener en cuenta que los archivos que se necesiten a la hora de correr los scripts estén en el mismo directorio que se vaya a ejecutar.



\apendice{Documentación de usuario}

\section{Introducción}
Este anexo va a describir la documentación referente a la ejecución de la aplicación desde el punto de vista del usuario.
\section{Requisitos de usuarios}
Los requisitos mínimos para poder usar la aplicación son:
\begin{itemize}
    \item Tener instalado Python 2.7
    \item Tener instaladas las bibliotecas Instagram-API-python, Networkx y Python-Louvain. En el anexo anterior se explica cómo se pueden conseguir dichas librerías.
    \item Tener conexión a internet.
    \item Tener una cuenta de Instagram con un mínimo de 5 personas seguidas que vayan a estar dentro de nuestra red.
    
\end{itemize}

\section{Instalación}
Para ejecutar la aplicación no es necesario cumplimentar una instalación tradicional, sino que con una descarga del repositorio de GitHub y ejecución de los distintos Scripts valdría.
\section{Manual del usuario}
Hay que cumplimentar unos pequeños pasos antes de ejecutar cada Script.
\subsection{NetCreator}
Requiere que se introduzca el usuario y contraseña de la cuenta de Instagram, así como las palabras claves que se quieren buscar para crear la red. Este primer paso lleva de varios días sí
\imagen{netcreator.png}{Login requerido}
\subsection{InsertFriends}
Requieres tanto el usuario y la contraseña de instagram como las que se refieren a MySQL, que es donde guardaremos la información que extraigamos.
\imagen{insertfriends.png}{Conexión a la base de datos}
\subsection{InsertFriendship}
Al igual que InsetFriends, necesitamos introducir el usuario y contraseña del Instagram y de MySQL.
Después de extraer la información de la API hay que configurar los datos como un CSV para que pueda leerlos bien Networkx. Para eso, se extraen desde la línea de comandos de la siguiente forma:
\begin{itemize}
    \item '/var/lib/mysql-files/prueba.csv' fields enclosed by '"' terminated by ';' escaped by '"' LINES terminated by '/r/n';
\end{itemize}


\subsection{Partition}
Partition es un script que no requiere de ningún cambio y que puede ejecutar sin problemas una vez hayamos cumplimentado todos los scripts anteriores.

\subsection{AnalizeCommunities}
Este es un notebook de Python que en vez de hacerlo un script como los anteriores, hay que subirlo a colab.research.google.com , puesto que necesita de mayores recursos que lo que normalmente se tiene en un ordenador personal. Nos requiere en el primer paso que le introduzcamos los csv de las particiones y al final nos devuelve los datos para descargar.

\subsection{Flask}
Para ejecutar Flask hay que escribir en la línea de comandos (desde Windows) hay que ir a la carpeta de ProyectoUBU dentro de flask y con botón derecho+mayusc abrir una powershell y escribir los 2 siguientes comandos:

\begin{itemize}
    \item set FLASK\_APP=flaskr
    \item flask run
\end{itemize}

Después abrir un navegador e ir a la página por defecto que es http://127.0.0.1:5000/
\imagen{flask.png}{Información mostrada al iniciar la web con Flask}

Se adjunta un video junto con los anexos que puede verse por si existieran dudas sobre la ejecución del programa.


\bibliographystyle{plain}
\bibliography{bibliografiaAnexos}

\end{document}
